\section{Experiments}
\label{sec:experiments}

\subsection{Datasets}

We assemble a heterogeneous corpus of paired videos
and 3D point trajectories spanning simulated and real-world
interactions.
The training corpus contains approximately 1.44 million
object-level samples, equivalent to 1,224 hours of
standardized 49-frame sequences at 16 FPS.
The data are organized into four collections:
\begin{itemize}
    \raggedright
    \item \textbf{Sim7} combines seven simulation datasets:
    ManiSkill~\citep{tao2024maniskill3},
    MIKASA~\citep{cherepanov2026memory},
    MetaWorld~\citep{yu2020meta},
    MetaFold~\citep{chen2025metafold},
    RLBench~\citep{james2020rlbench},
    RoboCasa~\citep{nasiriany2024robocasa}, and
    RoboTwin~\citep{mu2025robotwin},
    covering diverse robotic interactions and manipulation tasks.

    \item \textbf{Real3} comprises
    AgiBot World~\citep{bu2025agibot},
    EPIC-KITCHENS~\citep{damen2018epickitchens}, and
    Something-Something V2~\citep{goyal2017something},
    covering real-world robot and human interactions.

    \item \textbf{SpatialCog} is our self-collected dataset
    of tabletop and household interactions, with an emphasis
    on spatial relations in instruction-conditioned motion.

    \item \textbf{MolmoMotion-1M}~\citep{zhang2026molmomotion}
    contributes curated 3D motion annotations from six sources:
    EgoDex~\citep{hoque2025egodex},
    Xperience~\citep{ropedia2026xperience},
    YT-VIS~\citep{Yang_2019_ICCV},
    Stereo4D~\citep{Jin_2025_CVPR},
    DROID~\citep{khazatsky2024droid}, and
    MolmoSpaces~\citep{kim2026molmospaces}.
\end{itemize}

\paragraph{Annotation and standardization.}
For Sim7 and Real3, we generate captions with
Gemini-2.5 Pro~\citep{comanici2025gemini}, obtain object
masks with SAM 2~\citep{ravi2024sam2}, and extract 3D
point trajectories using Track4World~\citep{lu2026track4world}.
SpatialCog pairs manually annotated captions with
Track4World trajectories.
For MolmoMotion-1M, we reuse the released motion
annotations, apply source-specific geometry, scale,
and trajectory-validity checks, and standardize the
retained sequences into 49-frame clips.
All collections use a shared representation consisting
of video, a language instruction, initial 3D query points,
and frame-to-frame 3D displacements.
Additional processing details and per-source statistics
are provided in the supplementary material.

\paragraph{Evaluation split.}
We introduce WorldFlowBench, a benchmark comprising
1,678 held-out evaluation samples: 396 from Sim7,
882 from Real3, and 400 from SpatialCog.
The three subsets cover simulated interactions,
real-world motion, and spatial-relation reasoning,
respectively.
All evaluation samples are excluded from training.

\subsection{Training Configuration and Strategy}

\paragraph{Implementation details.}
We train on 49-frame clips at $288\times480$ resolution
with $N=64$ query points.
The video VAE uses temporal and spatial compression
factors of $q_t=4$ and $q_s=16$, followed by
$2\times2$ latent patchification ($p_s=2$).
Each clip therefore contains a $13\times9\times15$
video-token grid and $48\times64$ flow tokens.
Training uses 64 NVIDIA H800 GPUs across eight nodes,
with a per-GPU batch size of two and no gradient
accumulation, yielding a global batch size of 128.

\paragraph{Training strategy.}
We adopt a two-stage training schedule.
The first stage initializes the model from pretrained
Cosmos-Nano checkpoints and trains on the SpatialCog
training split for 32,000 steps.
The learning rate follows linear warmup and cosine decay,
with a peak of $10^{-4}$.
The resulting exponential moving average (EMA) weights
initialize the second stage.

The second stage trains on the heterogeneous data mixture
for 51,100 steps across three curriculum phases:
simulation-focused training, mixed real-world adaptation,
and refinement on the SpatialCog training split.
During real-world adaptation, sampling gradually shifts
from MolmoMotion to Real3.
Each phase prioritizes its target domain while retaining
cross-domain replay, with smooth transitions between
sampling distributions.
After warmup, the learning rate reaches $10^{-4}$,
transitions to an $8\times10^{-5}$ plateau during
real-world adaptation, and decays through late adaptation
and final refinement.

We expand task prompts into structured instructions
specifying data origin, camera motion, actor type,
and environment.
For classifier-free guidance, text conditioning is
dropped with probability $0.1$, while visual and
geometric conditioning is retained.
At inference, we use 40 rectified-flow sampling steps
with a guidance scale of $2$.



\subsection{Flow Benchmark}

\subsubsection{Evaluation Metrics}
We evaluate 3D point trajectories in metric coordinates.
Average displacement error (ADE) measures the mean
Euclidean position error across query points and future
frames, while final displacement error (FDE) measures
the mean position error at the final frame.
Both are reported in meters.
PWT measures the fraction of point--frame pairs whose
position error falls within a given distance threshold.
We report its average over thresholds of 1, 2, and 5~cm
to assess trajectory accuracy at fine spatial scales.

\subsubsection{Quantitative Evaluation on WorldFlowBench}
\label{sec:worldflowbench_eval}

We evaluate all learning-based methods under a single-frame
protocol, where one RGB observation and a language instruction
are provided to forecast a common future horizon without motion
history. All predictions are converted into metric 3D point
trajectories in the initial camera coordinate system. We
additionally report a Static reference that predicts zero flow
and therefore keeps every query point at its initial position
throughout the prediction horizon.

We compare against two families of learned baselines. The
video-generation-then-tracking family includes
NovaFlow~\citep{li2025novaflow}, as well as pipelines built on
Cosmos3-Nano~\citep{nvidia2026cosmos3} or
LingBot-Video-MoE~\citep{robbyant2026lingbotworld}, paired with
Track4World (T4W)~\citep{lu2026track4world} or
4RC~\citep{luo20264rc} for 3D tracking. The direct flow-generation
family includes MolmoMotion-AR-H1~\citep{zhang2026molmomotion},
$\mu_0$~\citep{lee2026mu_0}, and
EgoWAM~\citep{li2026egowam}.

As shown in Table~\ref{tab:worldflowbench}, \method{} achieves
the lowest ADE and FDE among all evaluated learning-based methods
on all three subsets. Relative to the strongest external learned
baseline on each subset, it reduces ADE by approximately
46.9\% on Sim7, 4.7\% on Real3, and 54.2\% on SpatialCog.
It also achieves the highest PWT among learning-based methods
on Sim7 and SpatialCog. On Real3, \method{} obtains a PWT of
0.311, compared with the best value of 0.332, while still
delivering the lowest ADE and FDE. These results demonstrate
consistent improvements in trajectory accuracy across simulated
interactions, real-world motion, and spatial-reasoning tasks.

\begin{figure}[t]
    \centering
    \includegraphics[width=\textwidth]{figures/qualitative_comparison.pdf}
    \caption{
        \textbf{Qualitative comparison on diverse real and simulated scenes.}
        We visualize predicted 3D query-point trajectories from $\mu_0$, MolmoMotion, and \method{} across SpatialCog, AgiBot, DROID, Stereo4D, Meta-World, RLBench, RoboTwin, and MetaFold.
        The trajectories are projected onto the corresponding input frames for visualization.
    }
    \label{fig:qualitative-comparison}
\end{figure}


\begin{table*}[t]
\centering
\small
\setlength{\tabcolsep}{4pt}
\renewcommand{\arraystretch}{1.12}
\resizebox{\textwidth}{!}{%
\begin{tabular}{@{}ll ccc ccc ccc@{}}
\toprule
\multirow{2}{*}{\textbf{Paradigm}} &
\multirow{2}{*}{\textbf{Model}} &
\multicolumn{3}{c}{\textbf{Sim7}} &
\multicolumn{3}{c}{\textbf{Real3}} &
\multicolumn{3}{c}{\textbf{SpatialCog}} \\
\cmidrule(lr){3-5}
\cmidrule(lr){6-8}
\cmidrule(lr){9-11}
& &
ADE$\downarrow$ & FDE$\downarrow$ & PWT$\uparrow$ &
ADE$\downarrow$ & FDE$\downarrow$ & PWT$\uparrow$ &
ADE$\downarrow$ & FDE$\downarrow$ & PWT$\uparrow$ \\
\midrule
Non-parametric
& Static
& 0.096 & 0.201 & 0.573
& 0.062 & 0.086 & 0.464
& 0.143 & 0.278 & 0.412 \\
\midrule
\multirow{5}{*}{\shortstack[l]{Video generation\\+ 3D tracking}}
& NovaFlow
& 0.196 & 0.318 & 0.289
& 0.095 & 0.129 & 0.293
& 0.165 & 0.295 & 0.320 \\
& Cosmos3-Nano + T4W
& 0.101 & 0.202 & 0.451
& 0.088 & 0.133 & 0.310
& 0.141 & 0.257 & 0.398 \\
& Cosmos3-Nano + 4RC
& 0.096 & 0.196 & \underline{0.531}
& 0.085 & 0.122 & \textbf{0.332}
& 0.141 & 0.263 & 0.401 \\
& LingBot-Video-MoE + T4W
& 0.105 & 0.207 & 0.425
& 0.097 & 0.145 & 0.267
& 0.134 & 0.239 & 0.364 \\
& LingBot-Video-MoE + 4RC
& 0.104 & 0.207 & 0.492
& 0.092 & 0.134 & 0.295
& 0.131 & 0.232 & 0.376 \\
\midrule
\multirow{3}{*}{\shortstack[l]{Direct flow\\generation}}
& MolmoMotion-AR-H1
& 0.176 & 0.323 & 0.280
& 0.185 & 0.335 & 0.222
& 0.282 & 0.533 & 0.173 \\
& $\mu_0$
& 0.186 & 0.307 & 0.069
& 0.116 & 0.156 & 0.148
& 0.213 & 0.270 & 0.055 \\
& EgoWAM
& 0.198 & 0.336 & 0.020
& 0.130 & 0.190 & 0.056
& 0.249 & 0.383 & 0.080 \\
\midrule
\multirow{4}{*}{Ablations}
& \method{} (flow-only)
& 0.099 & 0.206 & 0.495
& 0.101 & 0.152 & \underline{0.320}
& 0.136 & 0.248 & 0.372 \\
& \method{} + T4W
& 0.088 & 0.174 & 0.424
& 0.100 & 0.154 & 0.263
& 0.083 & 0.138 & 0.403 \\
& \method{} + 4RC
& 0.088 & 0.170 & 0.483
& 0.092 & 0.132 & 0.291
& 0.098 & 0.175 & 0.411 \\
& \method{} w/o auxiliary loss
& \underline{0.063} & \underline{0.116} & 0.527
& \underline{0.083} & \underline{0.115} & 0.299
& \underline{0.062} & \underline{0.084} & \underline{0.436} \\
\midrule
\textbf{Ours}
& \method{}
& \textbf{0.051} & \textbf{0.094} & \textbf{0.564}
& \textbf{0.081} & \textbf{0.113} & 0.311
& \textbf{0.060} & \textbf{0.078} & \textbf{0.449} \\
\bottomrule
\end{tabular}%
}
\caption{
\textbf{3D motion forecasting on WorldFlowBench.}
ADE and FDE measure average and final position errors
in meters, respectively; lower is better.
PWT is the mean point--frame hit rate over distance
thresholds of 1, 2, and 5~cm; higher is better.
Bold and underline mark the best and second-best results
in each column among non-static methods, including ablations.
Static is reported as a reference and excluded from ranking.
T4W denotes Track4World.
For \method{} + T4W/4RC, trajectories are extracted from
the videos generated by \method{} using the corresponding
tracker.
}
\label{tab:worldflowbench}
\end{table*}

\subsubsection{Ablation Studies}

\paragraph{Flow-only versus joint generation.}
To evaluate the benefit of jointly modeling 3D motion
and visual evolution, we compare \method{} with a
flow-only variant using the same backbone, training data,
and optimization recipe.
The variant predicts only 3D flow, whereas the full model
jointly generates video and flow.
The full model reduces ADE by 48.5\%, 19.8\%, and 55.9\%
on Sim7, Real3, and SpatialCog, respectively, with lower
FDE on all three subsets.
On SpatialCog, FDE decreases from $0.248$ to
$0.078\,\mathrm{m}$, a 68.5\% reduction.
PWT also improves on Sim7 and SpatialCog, with a small
decrease on Real3.
These results support the value of complementary visual
supervision for learning accurate 3D motion.

\paragraph{Direct flow generation versus post-hoc tracking.}
To evaluate flow as an explicit generative modality,
we compare directly generated trajectories with those
recovered by Track4World and 4RC from videos produced
by the same \method{} checkpoint.
This comparison holds the video generator and conditioning
inputs fixed, contrasting explicit motion generation
with motion recovery from generated pixels.
Direct flow generation achieves lower ADE and FDE
and higher PWT than both tracking pipelines on all
three subsets.
On SpatialCog, it achieves an ADE of $0.060\,\mathrm{m}$,
compared with $0.083\,\mathrm{m}$ for Track4World and
$0.098\,\mathrm{m}$ for 4RC.
These results support incorporating explicit 3D flow
into the world model's generative output rather than
recovering motion through a separate tracking stage.

\paragraph{Auxiliary motion supervision.}
We compare the full model with a variant trained without
auxiliary motion objectives.
These objectives supervise displacement, cumulative
motion direction, and trajectory endpoints through
the reconstructed clean flow.
Adding them improves ADE, FDE, and PWT on all three
subsets, with the largest ADE reduction on Sim7,
from $0.063$ to $0.051\,\mathrm{m}$ (19.0\%).
PWT increases from 0.527 to 0.564 on Sim7,
from 0.299 to 0.311 on Real3, and from 0.436 to 0.449
on SpatialCog.
These consistent improvements indicate that explicit
supervision of clean motion complements the primary
flow-matching objective.

% \subsubsection{PointMotionBench}

% We additionally evaluate \method{} on PointMotionBench
% under its established evaluation protocol.
% Table~\ref{tab:motion_prediction} includes the baseline
% groups and published results from MolmoMotion,
% alongside our evaluation of \method{}.
% Published results are distinguished from
% our own evaluations.

% \input{tables/pointmotionbench}

\subsection{Flow-Guided Robotic Manipulation}
\label{sec:real_robot}

\paragraph{Execution pipeline.}
We evaluate \method{} as an object-motion predictor within
a modular manipulation pipeline using a 7-DoF xArm,
an XHand dexterous hand, and a ZED Mini RGB-D camera.
Given an RGB-D observation and a language instruction,
we use Qwen3.8-27B~\citep{qwen38} to ground the target object and
obtain its mask. The mask and depth image are then used
to extract the target object's 3D point cloud, which
provides query points for \method{} to predict object
trajectories.
T(R,O) Grasp~\citep{fei2025t} generates the initial
dexterous grasp pose. Following
Dream2Flow~\citep{dharmarajan2025dream2flow}, a local
patch fitting optimizer fits a locally rigid region
near the grasp to the predicted object trajectories,
producing an arm trajectory that moves the grasped
object toward its predicted final position.
Thus, \method{} specifies object motion, while the grasp
and control modules translate it into robot execution.

\paragraph{Evaluation and results.}
We evaluate 12 tasks spanning pick and place, stacking,
and relative placement, with three randomized trials
per task and 36 executions per method. We evaluate flow
prediction and robot execution independently. Flow
prediction is considered successful if the predicted
object motion is reasonable for the given instruction,
while robot execution is considered successful if the
full manipulation pipeline completes the task on the
physical robot.
We select $\mu_0$~\citep{lee2026mu_0} as the baseline
based on its competitive offline trajectory prediction
performance among the reproducible direct flow prediction
methods evaluated in our benchmark. For each paired
trial, we replace only the flow predictor, keeping
initial layouts, instructions, grasp results, and
control settings identical.

As shown in Table~\ref{tab:real_robot}, \method{} achieves
a flow prediction success rate of 97.2\% (35/36),
compared with 25.0\% (9/36) for $\mu_0$. Its robot
execution success rate is 66.7\% (24/36), compared with
11.1\% (4/36) for $\mu_0$. The gap between flow
prediction and execution success for \method{} reflects
additional challenges in the manipulation pipeline,
including grasp failures, trajectory optimization
failures, and unstable placement during stacking.
These results show that \method{} provides more reliable
motion references and improves task completion within
the shared manipulation pipeline.

% Requires \usepackage{booktabs,graphicx}.
\begin{table*}[t]
\centering
\small
\setlength{\tabcolsep}{6pt}
\renewcommand{\arraystretch}{1.15}
\resizebox{\textwidth}{!}{%
\begin{tabular}{@{}lp{8cm}cccc@{}}
\toprule
\multirow{2}{*}{\textbf{Group}}
& \multirow{2}{*}{\textbf{Task Caption}}
& \multicolumn{2}{c}{\textbf{Flow Prediction}}
& \multicolumn{2}{c}{\textbf{Robot Execution}} \\
\cmidrule(lr){3-4}
\cmidrule(lr){5-6}
& & \method{} & $\mu_0$ & \method{} & $\mu_0$ \\
\midrule

Pick and Place
& Place the watermelon doll into the small pink plate.
& \textbf{3/3} & 0/3 & \textbf{3/3} & 0/3 \\
& Place the oolong tea into the large blue tray.
& \textbf{3/3} & 1/3 & \textbf{2/3} & 0/3 \\
& Place the potato chips into the dusty pink basket.
& \textbf{3/3} & 1/3 & \textbf{2/3} & 0/3 \\
& Place the yellow toy inside the red circle.
& \textbf{3/3} & 1/3 & \textbf{2/3} & 1/3 \\
\midrule

Stacking
& Place the green biscuit box on the translucent box.
& \textbf{3/3} & 1/3 & \textbf{3/3} & 0/3 \\
& Place the can of Red Bull on the translucent box.
& \textbf{3/3} & 1/3 & \textbf{2/3} & 1/3 \\
& Place the beef curry on the paper towel.
& \textbf{3/3} & 1/3 & \textbf{1/3} & 0/3 \\
& Place the yellow toy on the paper towel.
& \textbf{3/3} & 2/3 & 0/3 & 0/3 \\
\midrule

Relative Placement
& Place the sunscreen in front of the tissue.
& \textbf{2/3} & 0/3 & \textbf{1/3} & 0/3 \\
& Place the potato chips behind the Oreo box.
& \textbf{3/3} & 0/3 & \textbf{2/3} & 1/3 \\
& Place the watermelon doll to the left of the cola.
& \textbf{3/3} & 1/3 & \textbf{3/3} & 1/3 \\
& Place the oolong tea to the right of the small blue toy car.
& \textbf{3/3} & 0/3 & \textbf{3/3} & 0/3 \\
\midrule

\multicolumn{2}{@{}l}{\textbf{Overall success rate}}
& \textbf{97.2\%} & 25.0\%
& \textbf{66.7\%} & 11.1\% \\
\bottomrule
\end{tabular}%
}
\caption{
\textbf{Real-world manipulation with different flow predictors.}
Each task is evaluated over three trials, yielding 36 trials
per method. Flow prediction and robot execution are evaluated
separately for each trial. Both methods use matched initial
layouts, identical grasp results, and the same trajectory
optimization and execution settings; only the flow predictor
changes. Bold indicates the higher success rate.
}
\label{tab:real_robot}
\end{table*}